\section{2016: Obama's America}

I saw the surprising hit movie 2016: Obama's America this past weekend and am still unsure what its point is. It is the brainchild of the Anglo-Indian \textbf{Dinesh D'Souza} First of all, it was certainly a well-produced and researched documentary; as such, it is a model of what can—and should—be done by those who find themselves outside of Hollywood's consensus reality. There were several scenes focusing on the squalor of Indonesia and Kenya, hinting that they represent the future course of Obama's America.

First of all, we must point out that we reject the fundamental point of the film, \emph{viz}., that Obama's views are the result of some childhood psychological trauma. There is no fault in a boy's admiration for his father and there must be some compassion for boy's sense of loss for an absent father. Because of this, the film misses the opportunity to specifically point out the intellectual and ideological roots of Obama's worldview and how it differs from the neo-conservative viewpoint of D'Souza. Nor does it explain how members of Congress, not to mention millions of supporters, many with presumably happy childhoods, have also come to embrace Obama's vision for America.

Besides, although we could point out the psychological issues in D'Souza's childhood, there is something in D'Souza's worldview that colors his own perspective. He points out some accidental similarities between himself and Obama, such as their years of birth and marriage. However, we need to ``read between the lines", as it were, to get at the essential differences between them, as least how D'Souza sees them.

D'Souza was born into an Indian Catholic family in Goa, as his Portuguese name would imply. This made him neither fully an Indian nor a Westerner. He rejects the Indian traditions of arranged marriages and the caste system, neither of which applied to him, as well the insularity of village life. In India, a boy's life is mapped out from his birth, while in America, a boy can accomplish anything, depending only on his desires and talents. When he announced his imminent departure to America, his father told him he wouldn't like it: ``There are too many white people there." D'Souza fails to explain why the ``founders" whom he reveres, all but guaranteed that\footnote{\url{https://en.wikipedia.org/wiki/Naturalization_Act_of_1790}}. But that was then, this is now.

So we see that D'Souza's worldview is opposed to the alleged backwardness, or better said, the limitations, of so-called third-world countries and peoples. The USA, on the contrary, is based on an abstraction, the ``ideas" of the founding fathers; a man is an individual, not restricted to ties of family, caste, or religion, nor apparently even citizenship, as he got a position as a policy advisor to Ronald Reagan despite his status as an alien.

That may explain why D'Souza is so sensitive to anything in Obama that suggests the opposite. Specifically, the root of Obama's vision, according to D'Souza, is that Obama, on the contrary, valorizes the third world, to the extent even of threatening the economic dominance, cultural hegemony, and military might of the USA.

We can provide some quick examples. In foreign policy, for example, Obama allegedly favors Argentina in her dispute with the UK over the Falklands/Maldives. Clearly, this has no impact on American military security. Actually, in accordance with the long-standing Monroe doctrine\footnote{\url{https://en.wikipedia.org/wiki/Monroe_Doctrine}}, the USA should have opposed the introduction of British troops into the Falklands in the 1982 war, as Kennedy had done against the USSR in Cuba; at the very least, by that doctrine, the USA must remain neutral in the dispute.

From the unserious, he moves onto the serious. Putting up a map of countries with nuclear warheads, he shows Pakistan on one side and Israel, with 80 warheads, on the other side of Iran. For Iran, there is a question mark; actually, as of now, the count is zero. D'Souza asserts Obama is anti-Zionist, yet he doesn't make clear why the founders would have wanted the USA to become embroiled in that dispute.

The most serious charge is the Obama is deliberately trying to bankrupt the USA by incurring massive debt. This could be an intriguing film in itself, but D'Souza never explains why the Congress, including the Republicans, continues to permit that. It is true that eventually that debt will economically weaken the USA. Yet the policies that brought on the economic problems in the first place were allowed and promoted by the party that D'Souza supports.

D'Souza may, or may not, have made the case that Obama is unamerican and dangerous. But his strange brew of Neo-Christianity, eschatologically based Zionism, Randian Capitalism, and military adventurism may not be what is necessary to stop him.



\flrightit{Posted on 2012-09-07 by Cologero }

\begin{center}* * *\end{center}

\begin{footnotesize}\begin{sffamily}



\texttt{jwthomas on 2012-09-07 at 01:21 said: }

Neither Obama nor Romney have a ``worldview," unless it's gain power and do everything possible to stay in power. Their fuzzy opinions, like D'Souza's, are just intellectual constructions to justify and to cloak their real motives, which they rarely if ever admit to themselves. Is the ``hit movie" 2016 a hit with anyone other than those who already share its opinions?


\hfill

\texttt{Mihai on 2012-09-07 at 06:20 said: }

You said it in some other article.

This dispute (Democrats vs Republicans) is a dispute between the rising sudra domination vs the degenerate vaishya remnants.


\hfill

\texttt{Logres on 2012-09-07 at 10:38 said: }

Now that the movie is out, D'Souza is officially as mysterious as M. Night Shamalyan. I saw D'Souza debate someone back in the 90s on CSPAN, one of those ``liberal-conservative" cage matches – he made mincemeat of the guy, there is no doubt of his intelligence, \& although my dad loved this movie, I've avoided seeing it because of D'Souza's neo-con leanings on foreign policies. This is a good review.


\hfill

\texttt{Logres on 2012-09-07 at 10:41 said: }

D'Souza had a book out a few years ago that proposed that Muslims and Christians could cooperate against liberals, which got some people's attention.


\hfill

\texttt{Avery on 2012-09-07 at 12:02 said: }

D'Souza is probably quite literally a degenerated remnant of a vaishya family. But where would you place Mitt Romney in that scheme? He seems like a reconstructed vaishya, his spiritual strength coming from an undeveloped tradition…


\hfill

\texttt{Michael on 2012-09-07 at 12:57 said: }

Given that all modern politics, conservative or liberal, is really a variation on left-wing thought, shouldn't we still be involved in the political process to retain as many traditional elements as possible? After all, even Dante was involved in Florentine politics before he was banished.


\hfill

\texttt{Jason-Adam on 2012-09-08 at 22:23 said: }

It was right to engage in political activity when there was a chance of restoring Tradition, or at least preventing further collapse but now we've reach the end of the last age, all we can do is try to survive the coming collapse, and prepare to rebuild a new world when this one goes into ruin. Of course, I am talking as an American, to Americans. In other countries the ut look is not so bleak.


\end{sffamily}\end{footnotesize}
